%%%%%%%%%%%%%%%%%%%%%%%%%%%%%%%%%%%%%%%%%%%%%%%%%%%%%%%%%%%%%%%%%%%%%%
% Lớp 11 | Lý thuyết đồ thị | Đồ thị
% Dạng 1C2?1-1: Câu hỏi về đỉnh, cạnh của đồ thị
% Số câu: 4  (sinh tự động bởi tach_cau.py; từ nay thầy sửa trực tiếp file này)
%%%%%%%%%%%%%%%%%%%%%%%%%%%%%%%%%%%%%%%%%%%%%%%%%%%%%%%%%%%%%%%%%%%%%%

%%%%%%%%%%%%%%%% TỰ LUẬN %%%%%%%%%%%%%%%%

% nguồn: Ok_TuanLuan_Toan11_Lan1
\begin{bt}%[1C2B1-1]%[ma:4047978b]
	 \immini
	{ 
	Gọi $u_n$ là số hình tròn ở hàng thứ $n$ trong hình bên. Dự đoán công thức của số  hạng tổng quát cho dãy số $(u_n)$.
	} 
	{ 
	\begin{tikzpicture}[line join=round,line cap = round,>=stealth,scale=.8,font=\footnotesize,transform shape]
			 \foreach \i in {1,2,3,4}
				 {
			\draw (0,4.5 -  \i) node{$\text{Hàng thứ}\i$};
			\foreach \j  in {1,...,\i}
			 \draw[fill = green] (1 + \j,4.5 - \i) circle (7pt);
			} 
		\end{tikzpicture}
	 }
	 \loigiai{
		Ta thấy 
		\begin{itemize}
			\item Hàng $1$ có $1$ hình tròn. 
			\item Hàng $2$ có $2$  hình tròn.
			\item  Hàng $3$ có $3$ hình tròn.
			\item  Hàng $4$ có $4$ hình tròn.
		\end{itemize}
		Vậy ta dự đoán  công thức tổng quát của dãy số $(u_n)$ là $u_n = n$ với mọi $n\in\mathbb{N}^*$.
	 }
\end{bt}

% nguồn: Ok_TuanLuan_Toan11_Lan1
\begin{bt}%[1C2B1-1]%[ma:b53ab0dc]
	 \immini
	 {
	Gọi $v_n$ là tổng diện tích của các hình tô màu ở hàng thứ $n$ trong hình  bên (mỗi ô vuông nhỏ là một đơn vị diện tích). Dự đoán công thức của số hạng tổng quát cho dãy số $(v_n)$.
	 }
	{ 
	\begin{tikzpicture}[line join=round,line cap = round,>=stealth,scale=.8,font=\footnotesize,transform shape]
			\foreach  \i in {1,2,3,4}
				 {
			\draw (0,4.5 - \i) node{$\text{Hàng  thứ}\i$};
			 \pgfmathsetmacro\x{.1*\i}
			\foreach \j  in {1,...,\i}
				 {
			\begin{scope}[shift={(1 + \j,4.5 - \i - \x)}]
				\draw[fill = cyan!70,scale=.2] (0,0)  rectangle (\i,\i);
				\draw[scale=.2] (0,0) grid (\i,\i); 
			\end{scope}
			 }
			 }
		\end{tikzpicture}
	 }
	\loigiai{ 
		 Ta thấy
		\begin{itemize}
			\item Hàng $1$ có $1$ hình vuông cạnh $1$  đơn vị nên $v_1 = 1\times 1 = 1^3$.
			\item Hàng $2$ có $2$  hình vuông cạnh $2$ đơn vị nên $v_2 = 2\times 2^2 = 2^3$.
			\item Hàng $3$ có $3$ hình vuông cạnh $3$  đơn vị nên $v_3 = 3\times 3^2 = 3^3$.
			\item Hàng $4$ có  $4$ hình vuông cạnh $4$ đơn vị nên $v_4 = 4\times 4^2 = 4^3$.
		\end{itemize}
		Vậy ta dự đoán công thức tổng quát của dãy số $(v_n)$ là $v_n = n^3$  với mọi $n\in\mathbb{N}^*$.
	 }
\end{bt}

% nguồn: Ok_TuanLuan_Toan11_Lan2
\begin{ex}%[1C2B1-1]%[ma:9eb3e428]
	\immini 
	 {
	Gọi $u_n$ là số hình tròn ở hàng thứ $n$ trong hình  bên. Dự đoán công thức của số hạng tổng quát cho dãy số $(u_n)$.
	 }
	{ 
	\begin{tikzpicture}[line join=round,line cap = round,>=stealth,scale=.8,font=\footnotesize,transform shape]
			\foreach \i in  {1,2,3,4}
				{ 
			\draw (0,4.5  - \i) node{$\text{Hàng thứ}\i$};
			\foreach \j in  {1,...,\i}
			\draw[fill = green] (1 + \j,4.5 -  \i) circle (7pt);
			 }
		\end{tikzpicture}
	} 
	 \loigiai{
		Ta  thấy
		\begin{itemize}
			\item Hàng  $1$ có $1$ hình tròn.
			\item Hàng $2$ có $2$ hình  tròn.
			\item Hàng $3$ có $3$  hình tròn.
			\item Hàng $4$ có $4$ hình tròn. 
		\end{itemize}
		Vậy  ta dự đoán công thức tổng quát của dãy số $(u_n)$ là $u_n = n$ với mọi $n\in\mathbb{N}^*$.
	 }
\end{ex}

% nguồn: Ok_TuanLuan_Toan11_Lan2
\begin{ex}%[1C2B1-1]%[ma:7d5461af]
	 \immini
	{ 
	Gọi $v_n$ là tổng diện tích của các  hình tô màu ở hàng thứ $n$ trong hình bên (mỗi ô vuông nhỏ là một đơn vị diện tích). Dự đoán công thức của số hạng tổng quát cho dãy số $(v_n)$.
	 }
	 {
	\begin{tikzpicture}[line join=round,line cap = round,>=stealth,scale=.8,font=\footnotesize,transform shape]
			\foreach \i in {1,2,3,4} 
				{ 
			\draw (0,4.5 - \i) node{$\text{Hàng  thứ}\i$};
			 \pgfmathsetmacro\x{.1*\i}
			\foreach \j in {1,...,\i} 
				 {
			\begin{scope}[shift={(1 + \j,4.5 - \i - \x)}]
				\draw[fill =  cyan!70,scale=.2] (0,0) rectangle (\i,\i);
				\draw[scale=.2] (0,0) grid  (\i,\i);
			\end{scope}
			 }
			} 
		\end{tikzpicture}
	 }
	\loigiai{ 
		Ta thấy 
		\begin{itemize}
			\item Hàng $1$ có $1$ hình vuông cạnh $1$ đơn vị nên $v_1 = 1\times 1 = 1^3$. 
			\item  Hàng $2$ có $2$ hình vuông cạnh $2$ đơn vị nên $v_2 = 2\times 2^2 = 2^3$.
			\item Hàng $3$ có  $3$ hình vuông cạnh $3$ đơn vị nên $v_3 = 3\times 3^2 = 3^3$.
			\item Hàng $4$ có $4$ hình vuông cạnh $4$ đơn vị nên $v_4 =  4\times 4^2 = 4^3$.
		\end{itemize}
		Vậy ta dự đoán công thức tổng quát của dãy số $(v_n)$  là $v_n = n^3$ với mọi $n\in\mathbb{N}^*$.
	 }
\end{ex}
